\scanpage{5}
% source lines 1-2
Par suite, on a
\[
\varprojlim_{\underline m}{\unsure{\mathscr{P}|\mathfrak{P}|\mathcal{P}}}_K^{(m)}=C_K/\substack{\text{composante}\\\text{de l'unité}}\,D_K=C'_K.
\]

% source lines 3-9
Si $L$ est une extension \deleted{galoisienne} \inserted{finie} de $K$, on a
\[
\mathscr{P}_K^{\underline m}/N_{L/K}\mathscr{P}_L^{\underline m'}\approx C_K/N_{L/K}C_L
\]
pourvu que $N_{L/K}{\unsure{J|\mathscr{J}}}_L\supset U^{\underline m}$ (ce qui est vrai si $\underline m$ est assez grand, --- il suffit de le prendre grand \unsure{uniquement} aux \unsure{premiers sur} les places \uline{ramifiées} de $L$) et $\underline m$ étant \unsure{choisi ainsi}, $\underline m'$ (idéal absolu de $L$) est pris assez grand pour que $N_{L/K}L^{*\underline m'}\subset K^{*\underline m}$.

% source lines 10-15
\inserted{Supposons $L$ galoisienne.} On a un hom canonique (si $S$ désigne l'ens des places \inserted{\unsure{non}} ramifiées finies de $L$).
\[
{\unsure{D|\mathscr{D}}}^{\deleted{\unsure{\underline m}}S_0}\longrightarrow G(L/K)
\]
par la condition que si $\mathfrak p\notin S_0$, $\mathfrak p\unsure{-|\mapsto}(\mathfrak p,L/K)$ (substitution de Frobenius). Le noyau contient \deleted{\unsure{les}} $N_{L/K}D_{L\deleted{\unsure{K}}}^{S'_0}$ (trivial). La théorie du corps de classes

% source lines 16-20, circled a
\textcircled{a} affirme qu'on peut trouver $\underline m$ (\unsure{formé|premier} avec $S$) tel que \deleted{\unsure{tout} $(x)$} $P_K^{\underline m}$ est dans le noyau; elle implique
\diagram{\begin{tikzcd}[ampersand replacement=\&]
\text{qu'on a un hom. canonique}\&\mathscr{P}_K^{\underline m}\arrow[r]\arrow[d,phantom,"\wr" description]\&G(L/K)\\
\text{d'où un hom canonique}\&C_K/\operatorname{Im}U_K^{\underline m}\arrow[r]\&G(L/K)
\end{tikzcd}}

(et en fait un hom.
\diagram{\begin{tikzcd}[ampersand replacement=\&]
\mathscr{P}_K^{\underline m}/\unsure{N_{L/K}|\deleted{N_{L/K}}}\mathscr{P}_{\unsure{K|L}}^{m'}\arrow[r]\arrow[d,phantom,"\wr" description]\&G(L/K)\\
C_K/N_{L/K}C_L\arrow[ur]\&
\end{tikzcd}}
)

% source lines 21-26
% slanted
\margin{\uline{hom. de réciprocité}}
\unsure{Cela|Ceci} permet donc de définir, pour \deleted{\unsure{tout}} \uline{idèle} $\underline a$, \unsure{les éléments|l'élément} $(\underline a,L/K)\in G(L/K)$, \deleted{et} $(\underline a,L/K)_{\mathfrak p}=(\underline a_{\mathfrak p},L/K)$ (``\unsure{Normenrestsymbol}'' de Hilbert), c'est un hom. continu, défini par les conditions : $(\underline a_{\mathfrak p},L/K)=(\mathfrak p,L/K)$ si $\mathfrak p$ fini non ramifié \inserted{et} $v_{\mathfrak p}(a_{\mathfrak p})=1$, et $(\unsure{x},L/K)=\unsure{1|0}$ si \deleted{$\unsure{a}=$} $x\in K^{*\underline m}$ ($\underline m$ grand). De plus

% source lines 27-28, circled b
\textcircled{b} L'hom $D_K^{S_0}\longrightarrow G(L/K)$ est surjectif, (et même on a le ``\uline{th. de la progression arithmétique}'') et

% source lines 29-30, circled c
\textcircled{c} Son noyau est \uline{exactement} $P_K^{\underline m}N_{L/K}{\unsure{P|D}}_L^{m'}$, i.e. on a un \uline{isomorphisme}
\[
\mathscr{P}_K^{\underline m}/N_{L/K}\mathscr{P}_L^{m'}=C_K/N_{L/K}C_L\longrightarrow G(L/K).
\]
% archive sheet number: 4
